Consider a compact object (such as black hole, white dwarf or neutron star)
with a spherically symmetric accretion of gas, assume that accreting gas is
hydrogen. As particles fall into the object they heat up and radiate, thus
creating radiation pressure acting on rest of the accreting material. This
force is given by,
\[ F_L = \sigma_e \frac{I}{c} \]
where,
\[ \sigma_e = \frac{8\pi}{3}
   \left(\frac{e^2}{4\pi\epsilon_0 m_e c^2}\right)^2 \]
is the Thompson cross section for electrons, $c$ is the speed of light and $I$
intensity of the light. Although $F_L$ is calculated for electrons, it
effectively acts on the whole atom.

\begin{description}
\item[(a)] If the central compact object has a mass $M$, find an expression
  for the Eddington limit ($L_E$), which is the maximum possible luminosity
  for the accretion sphere. \hfill[3\,pt]
\item[(b)] For a particular compact object, the luminosity due to material
  accretion is the same as the Solar luminosity $L_\odot$. What is its minimum
  possible mass of this object to achieve this luminosity (in units of
  $M_\odot$)? \hfill[4\,pt]
\end{description}

Assume that the atoms in the accretion sphere originate far away from the
compact object. When these atoms fall into the compact object, their
gravitational energy is transformed into radiation.

\begin{description}
\item[(c)] Derive an expression for accretion luminosity ($L_{\mathrm{acc}}$)
  in terms of the compact object's mass ($M$), mass accretion rate
  ($\dot{M} = \frac{\Delta M}{\Delta t}$), and the compact object's radius
  ($R$). \hfill[2\,pt]
\item[(d)] Show that the maximum possible accretion rate in steady state is
  not directly dependent on the mass of the compact object. \hfill[2\,pt]
\item[(e)] For an object with $R = 12 \times 10^9$\,m, calculate the maximum
  possible accretion rate $\dot{M}$ (in units of $M_\odot$ per year).
  \hfill[2\,pt]
\end{description}

In reality, the accretion geometry is disk shaped, where most particles trace
an almost circular orbit around the compact object. Consider a binary system
consisting of a compact object of mass $M_1$ and a hydrogen burning star of
mass $M_2$ at a distance $a$ from each other. The gas from the hydrogen
burning star is accreted by the compact object and due to this mass transfer
the period of the binary system changes.

\begin{description}
\item[(f)] Find the condition on the two masses such that the separation
  between the stars is increasing. Ignore the rotation of the stars.
  \hfill[7\,pt]
\end{description}
